% a07_k01steps: proofs of the two statements of ssec:k01:open: the orbit regrouping of cut renormalised families
% (hyp:k01:RENorb) and the interaction-term half of Balaban's domain-change estimate at a nonzero background (hyp:k01:BALint).
% Part 1 of 4; the parts sec/a07_k01steps_a2.tex, sec/a07_k01steps_b.tex, sec/a07_k01steps_c.tex are input from here.
\providecommand{\BALint}{\ensuremath{(\mathrm{BAL}_{5})[\mathrm{int}]}}
\providecommand{\RENorb}{\ensuremath{(\mathrm{REN}_{\mathrm{orb}})}}
\providecommand{\Cl}{\operatorname{Cl}}
\providecommand{\cutg}{\operatorname{cut}}
\providecommand{\OmA}{\ensuremath{\Omega_{A}}}
\providecommand{\OmC}{\ensuremath{\Omega_{C}}}
\providecommand{\Bt}{\ensuremath{\tilde B}}
\providecommand{\dG}{\ensuremath{\delta G}}
\providecommand{\dE}{\ensuremath{\delta E}}
\providecommand{\Lsev}{\ensuremath{\Lambda^{(k)}_{7}}}

\section{Two steps of the renormalisation-group analysis}\label{app:a07_k01steps}

This appendix proves the two statements of \cref{ssec:k01:open}. Both concern Ba{\l}aban's renormalised classes of graphs \cite[Sects.~1--3]{Ba3} at the moment when some of their lines are \emph{cut}: replaced by a pair of external legs carrying a two-variable kernel. This happens when two propagator domains are compared (\cite[p.~414]{Ba3}: ``we will have also the propagators $\delta G_k(\Omega,\Omega_2,\tilde B)$'', which ``will be treated as an external field'') and whenever a line object is telescoped. \Cref{ssec:a07:orb} proves \RENorb\ (\cref{hyp:k01:RENorb}): the cut terms of a renormalised family regroup into complete renormalised classes of the cut graphs. \Cref{ssec:a07:dom} proves \BALint\ (\cref{hyp:k01:BALint}), the interaction-term half of \cite[Prop.~2.6]{Ba2} at a nonzero background, by an exact regrouping of a second kind, into full classes of the cut components, to which Ba{\l}aban's renormalisation bound \cite[Prop.~1, (1.33)]{Ba3} applies at every background.

\subsection{Orbit regrouping of cut renormalised families}\label{ssec:a07:orb}

\paragraph{Graphs and labels.} A \emph{renormalised family} is a finite sum $\cR=\sum_{\Gamma\in\cG}c_\Gamma E(\Gamma)$ over a finite set $\cG$ of graphs built from the vertices \cite[(1.6)--(1.15)]{Ba3} (the seagulls included), from mass-counterterm vertices \cite[(1.7)]{Ba3} with the coefficients \cite[(1.23), (1.29)]{Ba3}, and from source vertices; $c_\Gamma$ collects the couplings, the combinatorial weight and the counterterm values, and every counterterm is a vertex of the action with a field- and domain-independent coefficient (\cite[(1.29)--(1.31), p.~419]{Ba3}: the counterterms are built from torus propagators at zero background). Every vertex carries \emph{labels} \cite[p.~420]{Ba3}: each vector-field leg the support box (a grid box of side $\le2$) of its partition function, each vertex of type (1.6)/(1.7) its unit cube, each block vertex (1.13)--(1.15) its block $B^k(y)$; a vertex is identified by its type and its labels. The counterterm vertex that replaces a divergent two-scalar-leg subgraph $\gamma$ carries $\gamma$'s labels (\cite[p.~417]{Ba3}: ``the same graph $G$ but with both external legs localized in $x$'', and \cite[p.~420]{Ba3}: the localisation is performed ``for the terms of $\delta m^2_k$ also''), so that a graph and its counterterm partner have the same label multiset: this is Ba{\l}aban's ``localizations in these graphs should be in agreement'' \cite[p.~432]{Ba3}. A \emph{line} is a pair of legs with a kernel type (scalar or photon, with its domain and background); an unpaired leg is an \emph{external leg} and carries an external field, which may be a function of several variables \cite[p.~419]{Ba3}. The lines are full lines: the scale decomposition is internal to the bounds that are applied to the classes.

\paragraph{Subgraphs and degrees.} A \emph{subgraph} $H$ of $\Gamma$ is a connected set of lines of $\Gamma$ (a line set, not a vertex set); $V(H)$ is the set of its vertices, and its \emph{slots} are the legs of $V(H)$ that are not on lines of $H$: legs on other lines, external legs and cut ends. Its \emph{type} is the number and kind (scalar or photon) of its slots, with the differentiation of its scalar slots; its order is the sum of the orders of its vertices; its degree $D(H)$ is Ba{\l}aban's \cite[(2.1)--(2.3)]{Ba3}: $D_H(v)$ counts the factors $\eta$ of $v$, the legs of $v$ ``belonging to internal lines'' of $H$ (with weight $(2-d)/2$), minus the differentiations of $v$ acting on internal lines, plus the vector legs on internal lines when $v$ is a block vertex; ``external fields have a dimension equal to $0$''; $D(H)=\sum_{v}D_H(v)-d$; and a counterterm vertex counts as the graph it replaces. Hence \emph{$D(H)$ depends only on $(V(H),H,\text{which legs are slots})$}, and so does the type. The divergent subgraphs and their completions \cite[pp.~420, 429--432]{Ba3} are: a two-scalar-leg, one-particle-irreducible $H$ with undifferentiated slots and $D(H)\le0$, completed by its \emph{counterterm partner}; a two-vector-leg $H$ with $D(H)\le0$, completed by the \emph{two-vector-leg class}, in Ba{\l}aban's maximal form ``all possible connected graphs of the considered type \dots\ of a fixed order'' \cite[p.~431]{Ba3}; a $\phi\phi A$ subgraph with $D(H)=0$, one-particle-irreducible (\cite[(2.18)--(2.20)]{Ba3}) or reducible (\cite[(2.21)--(2.22)]{Ba3}), completed by the \emph{permutation} of its two external scalar legs \cite[p.~429]{Ba3}. Graphs with a block vertex have $D>0$ \cite[Cor.~2.3]{Ba3}, so no divergent subgraph contains one; vacuum subgraphs do not occur below (every member and every cut component carries an external field).

\paragraph{Moves.} Each completion is a \emph{move} $\mu$ attached to a subgraph $H$ of its source $\Gamma$: it replaces $H$ by a graph $H^*$ together with a bijection $s_\mu$ from the slots of $H$ to the slots of $H^*$, and whatever sits at a slot of $H$ (an outer line end, an external field, a cut end) is re-attached at the image slot; nothing else changes. (P) the counterterm partner: $V(H)$ and the lines of $H$ are replaced by one counterterm vertex carrying $H$'s labels, the two slots going to its two legs; (P$'$) its inverse, the expansion of a counterterm vertex; (W) the two-vector-leg completion: $H\to H'$ for every $H'$ in the per-label two-vector-leg class of $H$ (all connected graphs with $H$'s label multiset and order whose two free photon legs carry $H$'s two slot labels, with their counterterm members), the slots mapped by their vector-leg labels; it is applied at \emph{every} divergent two-vector-leg subgraph, maximal or not (substituting at a non-maximal one gives a member of the class of the maximal one, so the classes are the same); (X) the exchange: the attachments at the two scalar slots of a divergent $\phi\phi A$ subgraph are exchanged (if the two slots are joined by one line, that line is mapped to itself with its ends exchanged; if both carry external fields of one kind, their arguments are exchanged). Let $M$ be a set of moves. The moves (P), (W), (X) have the following five properties, which are all that the proofs below use: \textup{(A1)} a move changes only the lines and vertices of $H$ and the attachment of the slots, and does not depend on what is attached at the slots; \textup{(A2)} whether a move is available at $H$, and its results $(H^*,s_\mu)$, depend only on $H$ with its slots (lines, vertices, labels, slot types), because $D(H)$ and the slot types are intrinsic and the per-label two-vector-leg class depends only on $H$'s labels; \textup{(A3)} every move has an inverse in $M$ at $H^*$ with slot map $s_\mu^{-1}$; \textup{(A4)} a move keeps the label multiset, the kind of every slot and the kernel type of every surviving line; \textup{(A5)} $c_{\mu\Gamma}/c_\Gamma$ depends only on $(H,H^*,s_\mu)$ (the exchange keeps the labelled vertex multiset, and (P), (W) change it only inside $H$).

\begin{definition}[classes]\label{def:a07:class}
The \emph{class} $\Cl_M(\Gamma)$ of a graph is its connected component under the moves of $M$, the moves at $\Gamma$ itself included when $\Gamma$ is of a primitively divergent type (its one-vertex counterterm partner \cite[p.~430]{Ba3}, its permuted graph \cite[p.~429]{Ba3}, all two-vector-leg graphs of its order \cite[p.~431]{Ba3}). With $M$ the moves (P), (W), (X), $\Cl_M(\Gamma)$ is Ba{\l}aban's class $G_{\mathrm{ren}}$ of $\Gamma$: ``the smallest sets of graphs, symmetric with respect to permutation of external scalar field legs and with each divergent subgraph renormalized'' \cite[p.~429]{Ba3} are the connected components, since the moves are invertible. A \emph{per-label full class} is the set of all connected graphs of one type, one order and one label multiset, with their counterterm members \cite[p.~431]{Ba3}; it is closed under every move. A family is \emph{closed under $M$} if $\cG$ is a finite union of classes; Ba{\l}aban's renormalised families are.
\end{definition}

\begin{definition}[cut terms]\label{def:a07:cut}
Fix a set of \emph{decomposed kinds} of objects of the family (for instance all lines and legs on a domain, or all restricted lines), and write every object $o$ of such a kind as $o=o'-\Delta o$, with $\Delta o$ a symmetric kernel, $\Delta o(x,y)=\Delta o(y,x)^{T}$ (true for every difference of two symmetric resolvents and for the kernels used below). An \emph{oriented cut term} is $\tau=(\Gamma,S,\theta)$ with $S$ a set of decomposed objects of $\Gamma$ and $\theta$ an orientation of every line of $S$; its weight is $2^{-|S_L|}c_\Gamma$, $S_L$ the set of lines of $S$. By the symmetry of $\Delta o$ the two orientations of a line add up to $\Delta o$, so multilinearity gives the exact expansion
\begin{equation}\label{eq:a07:tel}
\cR\ =\ \sum_{\tau}2^{-|S_L|}c_\Gamma\prod_{o\in S}(-\Delta o)_\theta\prod_{o\notin S}o'.
\end{equation}
The \emph{cut graph} $\cutg(\tau)=\Gamma\setminus S$ is $\Gamma$ with the lines of $S$ removed, their ends carrying the two arguments of $\Delta o$ (end $0$ and end $1$ by $\theta$), each leg of $S$ carrying $\Delta o$, and the objects of $S$ labelled. This is Ba{\l}aban's ``we replace these lines by pairs of external legs'' \cite[p.~432]{Ba3}, with external fields ``in the form of functions of many variables'' \cite[p.~419]{Ba3}. A subgraph $H$ of $\Gamma$ \emph{avoids} $S$ if it contains no line of $S$. For a move $\mu$ at a subgraph avoiding $S$ put $\mu\cdot\tau\defeq(\mu\Gamma,s_\mu(S),s_\mu(\theta))$: the ends of $S$-lines and the $S$-legs sitting at slots of $H$ are moved to the image slots, everything else is fixed. The \emph{orbits} are the classes of the equivalence relation generated by $\tau\sim\mu\cdot\tau$.
\end{definition}

\begin{lemma}[orbit regrouping, \RENorb]\label{lem:a07:orb}
Let $\cR$ be closed under $M$, where $M$ satisfies \textup{(A1)--(A5)}. Then:
\begin{enumerate}[label=\textup{(\alph*)},nosep,leftmargin=*]
\item A connected line set $H$ of $\Gamma$ avoids $S$ if and only if it is a connected line set of $\cutg(\tau)$; its vertices, slots, slot types, order and degree are the same in $\Gamma$ and in $\cutg(\tau)$, and so are the moves available at it and their results.
\item $\mu\cdot\tau$ is an oriented cut term of $\cR$; $\mu^{-1}$ acts at $H^*$, which avoids $s_\mu(S)$, and $\mu^{-1}\cdot(\mu\cdot\tau)=\tau$. The orbits partition the cut terms; each is finite, and the number of $S$-objects of each kind is constant on it.
\item $\cutg(\mu\cdot\tau)=\mu(\cutg\tau)$. Hence $\cutg$ maps the orbit of $\tau$ onto $\Cl_M(\cutg\tau)$, and re-joining the labelled $S$-objects maps $\Cl_M(\cutg\tau)$ onto the orbit exactly $r_\tau$-to-one, where $r_\tau\ge1$ is the order of the group of relabellings of identical $S$-objects (same kind, same kernel, same orientation pattern) that map $\Cl_M(\cutg\tau)$ to itself; and $\Cl_M(\cutg\tau)=\prod_K\Cl_M(K)$ over the connected components $K$ of $\cutg\tau$. In particular the cut graphs of one orbit form exactly one class of the (possibly disconnected) cut graph.
\item At fixed labels, after the sum over the separable indices of the cut kernels (write $\Delta o(x_0,x_1)=\sum_a c_aF_a(x_0)H_a(x_1)^*$),
\[
\cR^{(O)}\ \defeq\ \sum_{\tau'\in O}2^{-|S_L|}c_{\Gamma'}\prod_{S'}(-\Delta o)\prod_{O_{\Gamma'}\setminus S'}o'\ =\ (-1)^{|S|}2^{-|S_L|}\,c_*\,r_O^{-1}\prod_KE\big(\Cl_M(K);\text{its fields}\big),
\]
where $E(\Cl_M(K);\text{fields})$ is the class sum of $K$ with the fields $F_a$, $H_a$, the $S$-legs' kernels and the original fields at its slots, and $c_*$ is a constant. (At fixed indices, when $r_O>1$ and two identical $S$-objects carry different indices, the left side is the $r_O$-symmetrisation of the right side.)
\item $\sum_S\cR^{(S)}=\sum_O\cR^{(O)}$: \eqref{eq:a07:tel} summed in another order.
\item All members of an orbit carry the same labels, and an end of an $S$-object moves only from a slot of $H$ to a slot of $H^*$: under \textup{(X)} to the other scalar slot, under \textup{(P)} onto the counterterm vertex (which carries $H$'s labels), under \textup{(W)} to the leg with the same vector-leg label. No component is created, merged or split, and the label set of each component is kept.
\item Each orbit lies in the cut terms of one class of $\cG$; the cut terms of one class with a fixed number of $S$-objects of each kind are a finite disjoint union of orbits.
\end{enumerate}
\end{lemma}

\begin{proof}
(a) ``$H$ avoids $S$'' says that $H$ is a set of lines of $\Gamma$ none of which is in $S$, that is, a set of lines of $\Gamma\setminus S$; connectedness and $V(H)$ are read from $H$ alone. The slots of $H$ in $\Gamma$ are the legs of $V(H)$ not on lines of $H$; in $\cutg(\tau)$ they are the same legs (a leg on an $S$-line in $\Gamma$ is an external leg of $\cutg(\tau)$ carrying an end of $\Delta o$; a leg on a non-$S$ line outside $H$ is on the same line in both; an external leg is the same). The kind and the differentiation of a slot are properties of the leg (on a cut end the derivative acts on the external field), and $D(H)$ depends only on $(V(H),H,\text{slots})$; so the type of $H$ is the same in both graphs, and by (A2) the available moves and their results coincide.

(b) By (A2) and the closure of $\cG$, $\mu\Gamma\in\cG$. By (A1) every object of $\Gamma$ outside $H$ survives in $\mu\Gamma$, with the same kernel by (A4); $S$ consists of objects outside $H$, so $s_\mu(S)$ is a set of decomposed objects of $\mu\Gamma$ of the same kinds and $s_\mu(\theta)$ orients its lines: $\mu\cdot\tau$ is a cut term of $\cR$. By (A3), $\mu^{-1}$ is available at $H^*$; $H^*$ avoids $s_\mu(S)$ because the $S$-objects of $\mu\cdot\tau$ are images of objects outside $H$, none internal to $H^*$; applying $\mu^{-1}$ undoes the slot map. So $\sim$ is generated by a symmetric relation on a finite set, and $s_\mu$ keeps the numbers and kinds of the $S$-objects.

(c) In $\cutg(\tau)$, $H$ is a subgraph with the same slots, by (a). The move $\mu$ applied to $\cutg(\tau)$ re-attaches at the image slots whatever sits at the slots of $H$, outer lines, external fields and cut ends alike (A1). In $\mu\cdot\tau$ the $S$-ends at the slots of $H$ are moved by the same $s_\mu$, and the other $S$-ends are untouched. Removing the $S$-lines of $\mu\cdot\tau$ therefore gives the same graph as moving the cut ends in $\cutg(\tau)$: $\cutg(\mu\cdot\tau)=\mu(\cutg\tau)$. Conversely every move of $\cutg(\tau)$ is at a subgraph $H$ of $\cutg(\tau)$, which avoids $S$ by (a), so it is the cut of an orbit move. Hence $\cutg$ maps the orbit onto the connected component of $\cutg\tau$ under $M$, which is $\Cl_M(\cutg\tau)$. A move acts at a connected $H$, which lies in one component $K$ of $\cutg\tau$, and $H^*$ is connected and attached at the same slots; so a move changes only $K$, keeps the number of components and the label set of each (A4). Moves in different components act on disjoint data and commute; hence $\Cl_M(\cutg\tau)=\prod_K\Cl_M(K)$, every combination being reached by moving one component at a time. Finally, re-joining the two ends of each labelled $S$-line and restoring each $S$-leg's field inverts $\cutg$ on labelled objects; two labelled cut graphs of the class give the same unlabelled cut term if and only if they differ by a permutation of identical $S$-objects; the set of such permutations preserving the class is a group, since relabelling commutes with every move by (A1); so every cut term of the orbit has exactly $r_\tau$ preimages.

(d) By (A5), $c_{\Gamma'}=c_\Gamma\prod_{\text{moves}}w(H,H^*,s)$, and the moves in different components multiply; in the label picture $c_{\Gamma'}=c_*\prod_Kc(K')$ with $c(K')$ the product of the vertex weights of $K'$ (repeated labels are handled by the ordered-partition count at a constant cost). For fixed labels, $E(\cutg\tau';\text{fields})=\prod_KE(K';\text{fields})$: positions are summed within the fixed label boxes, the external fields are one-variable fields at fixed separable indices, and the lines of different components are different. The separable index of an oriented $S$-line is kept by every move ($s_\mu$ moves ends and keeps the line). Hence $\sum_{\Cl_M(\cutg\tau)}cE=c_*\prod_KE(\Cl_M(K))$; the left side counts each cut term of the orbit $r_\tau$ times by (c), and the sign $(-1)^{|S|}$ and the weight $2^{-|S_L|}$ are common to the orbit. When $r_O>1$ and identical $S$-objects carry different indices, the relabellings permute the indices, which gives the symmetrisation; after the sum over the indices the identity holds as stated.

(e) The orbits partition the index set of the finite sum \eqref{eq:a07:tel}, by (b).

(f) (A4) gives the labels; the slot maps of (P), (W), (X) are as listed; by (c) no component is created, merged or split.

(g) Every move is at a divergent subgraph of a member of a class of $\cG$, so it is a move of that class; fixing the numbers of $S$-objects by kind is compatible with $\sim$ by (b).
\end{proof}

\begin{remark}\label{rem:a07:orb}
(i) The orientation device is needed exactly when a move maps a cut term to itself with the two ends of an $S$-line exchanged (both ends of one $S$-line at the two scalar slots of a $\phi\phi A$ subgraph): with unoriented terms the orbit is a class only at the level of the symmetric two-variable kernel. (ii) If $S$ contains no line, the cut graph is $\Gamma$ itself with its external fields changed, and the moves at $\Gamma$ itself (the permutation of a top-level $\phi\phi A$ graph, which moves a difference leg from one slot to the other; the one-vertex partner of a two-scalar-leg graph) are orbit moves. (iii) A move at a divergent subgraph that contains an $S$-line is not an orbit move; the cut graph needs no completion there, and the $S$-terms of $\Gamma$ and of its partners at that subgraph fall into different orbits, each complete. (iv) The per-object grouping (collecting, for a fixed $S$, the cut terms of all members of a class) is \emph{not} closed under the moves: with the pair \cite[(2.18)]{Ba3} inside a larger member, an external line at the (1.8) slot cut, the exchange at the pair moves the cut end to the seagull, so the per-object group contains the cut graph but not its partner, which is still needed, since the pair is still a divergent $\phi\phi A$ subgraph of the cut graph. Only the counterterm moves keep the per-object grouping, when the legs of a counterterm vertex are read as the legs of the graph it replaces. This is why the regrouping is by orbits.
\end{remark}

% a07_k01steps_a2: part of app:a07_k01steps (input from sec/a07_k01steps.tex): cut data of per-label full classes and the use of the two regrouping lemmas.

\begin{lemma}[cut data of a per-label full class]\label{lem:a07:full}
Let $\cG$ be a per-label full class of type $T$, order $o$ and label multiset $\cL$, and group the unoriented cut terms by their \emph{cut datum} $\delta$: the components $K$ of $\Gamma\setminus S$ with their label sets $\cL_K$ and their orders, at each $K$ the numbers of $S$-ends of each kind and of original external fields of each kind, and the incidence of the $S$-objects (which components each $S$-line joins, in which component each $S$-leg sits). Then:
\begin{enumerate}[label=\textup{(\alph*)},nosep,leftmargin=*]
\item Label the $S$-objects and fix one matching of their ends. The cut terms with datum $\delta$ correspond, $s_\delta$-to-one, to the tuples $\prod_K\cK_K$, where $\cK_K$ is the per-label full class of the component (all connected graphs with label multiset $\cL_K$ and the order of $K$ whose unpaired legs carry the prescribed fields, $S$-end kernels and original fields, in all assignments, with their counterterm members, \emph{members without internal line and of order $0$ included}), and $s_\delta$ is the order of the group of relabellings of identical $S$-objects of $\delta$, the exchange of the two ends of an $S$-line with both ends in one component included.
\item At fixed labels and separable indices, $\cR^{(\delta)}=(-1)^{|S|}c_*s_\delta^{-1}\prod_KE(\cK_K;\text{fields})$.
\item Every orbit of \cref{lem:a07:orb} lies in one group; every group is closed under the exchange of the two scalar legs of every two-scalar-leg or $\phi\phi A$ subgraph, the component itself included, and under the completion of a one-vector-leg subgraph to the two-vector-leg class of its first background-Taylor coefficient. Hence each group is a union of orbits also for the moves enlarged by these two.
\item $\sum_S\cR^{(S)}=\sum_\delta\cR^{(\delta)}$.
\end{enumerate}
\end{lemma}

\begin{proof}
(a) Cutting lands in $\prod_K\cK_K$: each component is connected, has the label set $\cL_K$ and the prescribed unpaired legs, and every counterterm vertex of $\Gamma$ lies in one component with labels inside $\cL_K$. Conversely, given members $K_i\in\cK_K$, an assignment of the labelled stub fields to their unpaired legs (included in $\cK_K$) and the fixed matching, gluing the $S$-lines gives a graph with label multiset $\bigcup_K\cL_K=\cL$, the uncut type $T$ and the order $o$, connected because the incidence multigraph of $\delta$ is connected (the original $\Gamma$ was); it is therefore a member of the full class $\cG$, with $S$ among its lines and legs. A two-scalar-leg subgraph of $K_i$ avoids $S$, so by \cref{lem:a07:orb}(a) it is a two-scalar-leg subgraph of $\Gamma$, and its counterterm member corresponds to $\Gamma$'s; conversely a counterterm vertex of $\Gamma$ is a vertex of one component. Two labelled tuples give the same unlabelled cut term if and only if they differ by a relabelling of identical $S$-objects, and such a relabelling maps the product of full classes to itself (each $\cK_K$ contains all assignments of its stub fields); so the map is $s_\delta$-to-one. (b) As in \cref{lem:a07:orb}(d). (c) The moves keep $\delta$: labels by (A4), stubs by \cref{lem:a07:orb}(f), incidence because $S$-lines survive. A full class contains every pairing of its legs, and the two further moves are re-pairings inside a component with the same labels and stubs (the exchange is a bijection of Wick pairings; the first background-Taylor coefficient of the full connected generating function of a one-vector-leg type is the full two-vector-leg class, since the generating function depends on the background and on the other vector legs only through their sum). (d) The groups partition the unoriented cut terms.
\end{proof}

\begin{remark}[how the two lemmas are used]\label{rem:a07:use}
When a renormalisation bound (\cite[Prop.~1, (1.33)]{Ba3} or its cube form) is applied to a telescoped family, it is applied per orbit: by \cref{lem:a07:orb}(c), (d) each orbit is a product of complete classes of the components, each carrying at least one external field (an $S$-end or an original field), with the cubes of the original graph, so the bound applies to each component with its tree decay, and the components are joined along a spanning tree of the cut graph's component multigraph by the decay of the cut kernels (\cite[(3.2), p.~432]{Ba3}: ``the exponential factor in (1.33) is obtained from the estimates of external fields''). The gain of a cut kernel is read at the maximum over the end cubes of the orbit, which is legitimate because the kernel bounds hold at every pair of cubes at the required depth (\cref{lem:a07:K3}); where a consumer needs the gain at one member's cubes, any two cubes of the fixed vertex set $W$ are within $d(W)$, so the conversion costs $e^{2|S|\delta'_\sigma d(W)}$ with the gain taken at the reduced rate $\delta'_\sigma=\delta'/(4|S|)$, which half of the tree factor $e^{-\delta'd(W)}$ absorbs. At zero background the orbits of (P), (W), (X) suffice. At a nonzero regular background the renormalisation bound is proved for per-label full classes (the touched two-scalar-leg subgraph becomes a $\phi\phi A$ node whose partner is its mirror, a move not among (P), (W), (X)); there the regrouping is by the groups of \cref{lem:a07:full}, and the members of the component classes without internal line or of order $0$ are treated as in the proof of \cref{lem:a07:S}.
\end{remark}


% a07_k01steps_b: part of app:a07_k01steps (input from sec/a07_k01steps.tex): the setting of Balaban's step, the exact background move, the elementary lemmas and the difference-kernel bound.

\subsection{The domain-change estimate at a nonzero background}\label{ssec:a07:dom}

\paragraph{The setting.} We are at the step $k\ge1$ of \cite[Sect.~2]{Ba2}, $d=3$, $L=5$, $t\defeq5^k\eps\le\eps_0$, $\eta\defeq5^{-k}$, in scale-$k$ units (unit blocks on $\eta\Z^3$; $\dist$ is the sup-metric \cite[(1.3)]{Ba1}, and $|\cdot|$ counts unit blocks). The couplings are $e_k=et^{1/2}$, $\lambda_k=\lambda t$; $p=p(t)=b_0(1+\ln t^{-1})^{\tilde p}$ is the small-field threshold and $r=r(t)=R(1+\ln t^{-1})^{\varrho}$, $\varrho>1$, $R>R_0$, the layer width \cite[(2.7), p.~558]{Ba2}; $R$ may be enlarged, and $r\ge R$. The large-field regions $\Lambda^{(k)}_0\supset\Lambda^{(k)}_1\supset\dots\supset\Lambda^{(k)}_7$ are those of \cite[(2.8), pp.~558, 566]{Ba2}: $(\Lambda^{(k)}_{i+1})^c$ is the union of the large blocks within distance $r$ of $(\Lambda^{(k)}_i)^c$; $\Lambda^{(k)}_0\subset\Lambda^{(k-1)\prime}_7$, primed sets being the scale-$(k-1)$ sets read at scale $k$; $B^k(S)$ is the set of fine-lattice points of the blocks of $S$. The two propagator domains of \cite[Prop.~2.6]{Ba2} are
\[
\OmA\defeq B^k(\Lambda^{(k-1)\prime}_2)\ \supset\ \OmC\defeq B^k(\Lambda^{(k)}_2),
\]
the inclusion because $\Lambda^{(k)}_2\subset\Lambda^{(k)}_0\subset\Lambda^{(k-1)\prime}_7\subset\Lambda^{(k-1)\prime}_2$; both are unions of large blocks. The background on $B^k(\Lambda^{(k)}_0)$ is $\Bt=(1-\theta_{k+1})A^{(k)}+\theta_{k+1}B^{(k+1),\eta}$ \cite[p.~579]{Ba2}, where $\theta_{k+1}$ equals $1$ on $B^k(\Lambda^{(k)}_2)$ and varies from $1$ to $0$ on a slice of thickness $<M$ around it (\cite[p.~567]{Ba2} at the step $k$, \cite[(2.96)]{Ba2} at the step $k+1$); $\Bt$ is regular on $\OmA$ in the sense of \cite[(2.23)]{Ba1}, with $|\partial\Bt|\le O(p)$ (\cite[(2.60)]{Ba2} for $A^{(k)}$ on $\OmA$, \cite[(2.94)--(2.98)]{Ba2} for $\Bt$ on $B^k(\Lambda^{(k)}_0)$). For a union $\Omega$ of large blocks and a regular field $A$ on it, $G_k(\Omega,A)=(-\Delta^{N}_{A,\Omega}+m^2t^2+a_kP_k(A))^{-1}$ is the scalar propagator of the step on $\Omega$ with Neumann conditions \cite[(2.22)]{Ba1}, \cite[(1.3)]{Ba4}, $P_k(A)$ the projection onto the transported block constants. For a function $F$ of one or two lattice variables on a set $S$, $\norm F_{1,\alpha;S}$ is the natural two-variable extension of Ba{\l}aban's norm \cite[(1.32), p.~420]{Ba3}: the maximum, over the variables, of the sup of $F$, of $|D_AF|$ in each variable, of $|D_xD_{x'}F|$, and of the transported H\"older quotients of exponent $\alpha$ of these derivatives in each variable; for a localisation $h$ ($h=1$ on a cube, smooth, supported in a fixed neighbourhood), $\norm{hF}_{1,\alpha}\le C_h\norm F_{1,\alpha;\supp h}$. We write $\sigma,\sigma'\in\{0,D,HD\}$ for ``no derivative'', ``the covariant derivative'' and ``the H\"older quotient of the covariant derivative'', with $|0|=0$, $|D|=1$, $|HD|=1+\alpha$.

\paragraph{Standing lower bounds.} $R\ge R_0+3$ and
\begin{equation}\label{eq:a07:D7}
D_7\ \defeq\ 5(r-1)-4\ \ge\ \max(R_0+2,\,r),
\end{equation}
where $R_0\ge1$ is the depth constant of \cite[Prop.~2.1]{Ba1}; since $r\ge R$ these are one more lower bound on $R$ or, equivalently, one more cap on $\eps_0$ of Ba{\l}aban's kind.

\paragraph{Printed inputs.} \textup{(P1)} \cite[Prop.~2.1, (2.24)--(2.26)]{Ba1} (which is the theorem \cite[(1.9)--(1.12)]{Ba4}): for a regular $A$ on $\Omega$ and a bounded source $f$, bounds on $G_k(\Omega,A)f$, $DG_k(\Omega,A)f$ and the H\"older quotient of the latter, with decay $c_0e^{-\delta_0\dist(x,\supp f)}$, at points at distance $\ge R_0$ from $\Omega^c$ (2.24)--(2.25), and for $\Omega_2\subset\Omega$ the difference $(G_k(\Omega,A)-G_k(\Omega_2,A))f$ with the additional factor $e^{-\delta_0\dist(\supp f,\Omega_2^c)-\delta_0\dist(\{x,x'\},\Omega_2^c)}$ (2.26), $\delta_0$, $c_0$, $R_0$ independent of $A$, $k$ and $\Omega$; these hold at every scale $j\le k$ after rescaling, by the same statement. \textup{(P2)} \cite[Props.~2.2--2.3]{Ba1} and \cite[(1.13)--(1.20)]{Ba4}: the exponential decay of the kernel of the block-spin operator $\Delta^{(j)}(\Omega,A)$ and of its domain difference $\Delta^{(j)}(\Omega,A)-\Delta^{(j)}(\Omega_0,A)$, the latter with the two-sided factor $e^{-\delta_0(\dist(x,\Omega^{(j)c})+\dist(x',\Omega^{(j)c}))}$ \cite[(2.27)--(2.29)]{Ba1}; the positivity and exponential decay of the restricted block covariances $C^{(j)}_\Lambda(\Omega,A)=\big((\Delta^{(j)}(\Omega,A)+aL^{-2}P_L)|_\Lambda\big)^{-1}$ with the fixed bare $a$ \cite[(2.31)--(2.34)]{Ba1}, \cite[(1.13), (1.15)--(1.16)]{Ba4}; and the two-sided decay of their domain differences $C^{(j)}_\Lambda(\Omega,A)-C^{(j)}_\Lambda(\Omega_0,A)$ \cite[(1.19)--(1.20)]{Ba4}; $\Lambda$ a union of big blocks, constants $c_2$, $\theta_0$ free of $j$, $\Omega$, $\Omega_0$, $\Lambda$ and $A$. \textup{(P3)} the Combes--Thomas bound of \cref{lem:a07:CT} below for the fixed-spacing resolvent, with no restriction near the boundary; the positivity it needs is \cite[(1.8)]{Ba4}. Below, $\theta_0$ denotes the smaller of the rates of \textup{(P2)} and \textup{(P3)}.

\begin{lemma}[a Combes--Thomas bound at fixed spacing]\label{lem:a07:CT}
Let $\Omega$ be a subset of a lattice of spacing $\eta\le1$ with the sup-metric, and on $\ell^2(\Omega)$ (weights $\eta^d$) let $T=D_A^*D_A+V+\sum_Bc_B\Pi_B$, where $D_A$ is the covariant difference on the bonds of $\Omega$, $V\ge0$ is a multiplication operator, the $\Pi_B$ are orthogonal projections whose ranges are supported in pairwise disjoint sets $B$ of diameter $\le\ell_*$, $\eta\le\ell_*$, $0\le c_B\le\bar a$, and $T\ge\gamma>0$. Put $\theta_0\defeq\min\{\ell_*^{-1},\gamma^{1/2}/(36\sqrt d),\gamma/(12e^2\bar a\ell_*)\}$. Then for all $f,f'\in\ell^2(\Omega)$ and all bond functions $g,g'$, with $\delta$ the distance of the supports (a bond located at its initial point),
\begin{align*}
|\ip f{T^{-1}f'}|&\le2\gamma^{-1}e^{-\theta_0\delta}\norm f\norm{f'},\\
|\ip g{D_AT^{-1}f'}|,\ |\ip f{T^{-1}D_A^*g'}|&\le(2\gamma^{-1/2}+6\sqrt d\,\theta_0\gamma^{-1})\,e^{\theta_0}e^{-\theta_0\delta}\,\norm\cdot\norm\cdot,\\
|\ip g{D_AT^{-1}D_A^*g'}|&\le4e^{\theta_0}e^{-\theta_0\delta}\norm g\norm{g'},
\end{align*}
with no restriction on the position of the supports relative to $\partial\Omega$ (the factor $e^{\theta_0}$ covers locating a bond at either end). In particular $\norm{\mathbf 1_YT^{-1}f}\le2\gamma^{-1}e^{-\theta_0\dist(Y,\supp f)}\norm f$ and $\norm{\mathbf 1_YD_AT^{-1}f}\le(2\gamma^{-1/2}+6\sqrt d\,\theta_0\gamma^{-1})e^{\theta_0}e^{-\theta_0\dist(Y,\supp f)}\norm f$ for every set $Y$. For the operators of \cref{lem:a07:K3} in scale-$1$ units ($\eta=1/L$, unit blocks, $\ell_*=1$, $c_B=a$, $V=m^2t^2l_1^2$) the lower bound $T\ge\gamma$ is \cite[(1.8)]{Ba4}, with $\gamma$ free of $\eta$, $\Omega$ and $A$.
\end{lemma}
\begin{proof}
\emph{Weight.} Let $\rho\defeq\dist(\cdot,\supp f')$ on the sites, $\rho(b)\defeq\rho(b_-)$ on the bonds, and $T_\theta\defeq e^{\theta\rho}Te^{-\theta\rho}$; then $T^{-1}=e^{-\theta\rho}T_\theta^{-1}e^{\theta\rho}$ and $e^{\theta\rho}f'=f'$. \emph{Hopping.} The ends of a bond are at distance $\eta$ and $\rho$ is $1$-Lipschitz, so for $\theta\eta\le1$, $e^{\theta\rho}D_Ae^{-\theta\rho}=D_A+\theta B_1$ with $(B_1u)(b)=\theta^{-1}\eta^{-1}\big(e^{\theta(\rho(b_-)-\rho(b_+))}-1\big)U_bu(b_+)$; since $|\theta^{-1}\eta^{-1}(e^{\theta s}-1)|\le e\le3$ for $|s|\le\eta$ and each site is the final point of $d$ bonds, $\norm{B_1}\le3\sqrt d$. Likewise $e^{\theta\rho}D_A^*e^{-\theta\rho}=(D_A-\theta B_1')^*$ with $\norm{B_1'}\le3\sqrt d$. \emph{The hopping change.} $e^{\theta\rho}D_A^*D_Ae^{-\theta\rho}-D_A^*D_A=\theta(D_A^*B_1-B_1'^*D_A)-\theta^2B_1'^*B_1$, so $|\ip\phi{(\cdots)\psi}|\le3\sqrt d\,\theta(\norm{D_A\phi}\norm\psi+\norm\phi\norm{D_A\psi})+9d\theta^2\norm\phi\norm\psi$. \emph{Pins.} Write $\Pi_B=WW^*$ and let $\rho_B$ be the value of $\rho$ at a point of $B$; then $e^{\theta\rho}\Pi_Be^{-\theta\rho}=(e^{\theta(\rho-\rho_B)}W)(e^{-\theta(\rho-\rho_B)}W)^*$, and since $|\rho-\rho_B|\le\ell_*$ on $B$ this differs from $\Pi_B$ by at most $(2e^x+xe^{2x})x\le2e^2x$ in norm, $x\defeq\theta\ell_*\le1$; the $B$ being disjoint, the change of the pin term has norm $\le2e^2\bar a\,\theta\ell_*$. \emph{Perturbation.} Put $K\defeq T^{-1/2}(T_\theta-T)T^{-1/2}$. Since $D_A^*D_A\le T$, $\norm{D_AT^{-1/2}}\le1$, and $\norm{T^{-1/2}}\le\gamma^{-1/2}$; hence $\norm K\le6\sqrt d\,\theta\gamma^{-1/2}+9d\theta^2\gamma^{-1}+2e^2\bar a\theta\ell_*\gamma^{-1}\le\frac16+\frac1{144}+\frac16<\frac12$ for $\theta\le\theta_0$. \emph{Bounded inverse.} $T_\theta=T^{1/2}(1+K)T^{1/2}$, so $T_\theta^{-1}=T^{-1/2}(1+K)^{-1}T^{-1/2}$ with $\norm{T_\theta^{-1}}\le2\gamma^{-1}$, $\norm{D_AT_\theta^{-1}},\norm{T_\theta^{-1}D_A^*}\le2\gamma^{-1/2}$ and $\norm{D_AT_\theta^{-1}D_A^*}\le2$. \emph{Bounds.} $\ip f{T^{-1}f'}=\ip{e^{-\theta\rho}f}{T_\theta^{-1}f'}$ with $\norm{e^{-\theta\rho}f}\le e^{-\theta\delta}\norm f$; $\ip g{D_AT^{-1}f'}=\ip{e^{-\theta\rho}g}{(D_A+\theta B_1)T_\theta^{-1}f'}$ with $\norm{(D_A+\theta B_1)T_\theta^{-1}}\le2\gamma^{-1/2}+6\sqrt d\,\theta\gamma^{-1}$, and the adjoint form follows from the symmetry of $T$; for the third form take $\rho\defeq\dist(\cdot,\supp g')$, so that $\ip g{D_AT^{-1}D_A^*g'}=\ip{e^{-\theta\rho}g}{(D_A+\theta B_1)T_\theta^{-1}(D_A-\theta B_1')^*g'}$ and $\norm{(D_A+\theta B_1)T_\theta^{-1}(D_A-\theta B_1')^*}\le2+12\sqrt d\,\theta\gamma^{-1/2}+18d\theta^2\gamma^{-1}\le4$. Take $\theta=\theta_0$.
\end{proof}

\begin{lemma}[the background move is exact]\label{lem:a07:C51}
On every bond of $\OmC$, $\Bt=B^{(k+1),\eta}$. Hence $G_k(\OmC,\Bt)=G_k(\OmC,B^{(k+1),\eta})$, $\Delta^{(k)}(\OmC,\Bt)=\Delta^{(k)}(\OmC,B^{(k+1),\eta})$ and $C^{(k)}_\Lambda(\OmC,\Bt)=C^{(k)}_\Lambda(\OmC,B^{(k+1),\eta})$ for every $\Lambda\subset\OmC^{(k)}$.
\end{lemma}
\begin{proof}
$\theta_{k+1}=1$ on $\OmC=B^k(\Lambda^{(k)}_2)$, and $\OmC\subset B^k(\Lambda^{(k)}_0)$, so $\Bt=B^{(k+1),\eta}$ there. The three operators depend on the field only through the transporters on the bonds of their domain (the Neumann Laplacian, \cite[(1.3)]{Ba4}) and on the contours $\Gamma^{(k)}_{y,x}$, which lie in $B^k(y)\subset\OmC$ \cite[(1.4)]{Ba4}.
\end{proof}
So the replacement of \cite[Prop.~2.6]{Ba2} is a pure domain change: $G_k(\OmA,\Bt)\to G_k(\OmC,\Bt)$, with $\dG\defeq G_k(\OmA,\Bt)-G_k(\OmC,\Bt)$ on $\OmC\times\OmC$ the $\delta G_k(\Omega,\Omega_2,\Bt)$ of \cite[p.~414]{Ba3}. The terms containing $\tilde B^{(k)}$ were estimated before it (\cite[p.~579]{Ba2}: ``Proposition I.3.1 can also be applied \dots\ the terms containing the configuration $\tilde B^{(k)}$ \dots\ by $O((L^k\eps)^\kappa)|\Lambda^{(k)}_2|$''), as were the terms of order $>\bar n$ and the terms localised in $B^k((\Lambda^{(k)}_7)^c)$; they are not part of the proposition.

\begin{lemma}[three elementary facts]\label{lem:a07:G}
\begin{enumerate}[label=\textup{(\alph*)},nosep,leftmargin=*]
\item \textup{(layer distances)} For $i\ge0$, $j\ge1$: $\dist(\Lambda^{(k)}_{i+j},(\Lambda^{(k)}_i)^c)\ge j(r-1)$; $\dist(\Lambda^{(k-1)\prime}_7,(\Lambda^{(k-1)\prime}_2)^c)\ge r-1$; and $\dist(B^k(S),B^k(S'))\ge\dist(S,S')-1$ on $\eta\Z^3$. In particular every point within distance $2+\tfrac14+\eta<3$ of $B^k(\Lsev)$ lies at depth $\ge D_7$ in $\OmC$.

\item \textup{(smallness)} For $\delta>0$, $n,m\ge0$ and $\kappa>0$ there is $t_0=t_0(\delta,n,m,\kappa,R,\varrho,b_0,\tilde p,\lambda)$ with $p(t)^n\lambda(t)^{-m}e^{-\delta r(t)}\le t^\kappa$ for $t\le t_0$; and for every $\beta>0$, $p(t)^n\le C_{\beta,n}t^{-\beta}$ for $t\le1$.
\item \textup{(tree sums)} Let $d(\{\square\})$ be the length of a shortest tree connecting given cubes of side $\le2$ \cite[p.~420]{Ba3}. For a fixed $\square_1$, $\sum_{\square_2,\dots,\square_n}e^{-\delta d(\{\square_1,\dots,\square_n\})}\le C_T(n,\delta)$, and for a set $S$ of unit cubes $\sum_{\square}e^{-\delta\dist(\square,S)}\le C(\delta)|S|$.
\end{enumerate}
\end{lemma}
\begin{proof}
(a) If $\dist(w,\Lambda_i^c)\le r$, the large block of $w$ lies within $r$ of $\Lambda_i^c$, so the closed $r$-neighbourhood of $\Lambda_i^c$ lies in $\Lambda_{i+1}^c$. Induction on $j$: for $x\in\Lambda_{i+j+1}$ and $z\in\Lambda_i^c$ nearest to $x$, the point $w$ on a sup-metric geodesic from $z$ to $x$ with $\dist(z,w)=\lfloor r\rfloor\ge r-1$ lies in $\Lambda_{i+1}^c$, so $\dist(x,z)\ge j(r-1)+(r-1)$. The second claim is the first at scale $k-1$ ($5(r(t/5)-1)$, divided by $5$, and $r(t/5)\ge r(t)$). The third holds because a block has sup-diameter $\le1$. For the last claim, the first claim with $i=2$, $j=5$ and the third give $\dist(B^k(\Lsev),\OmC^c)\ge5(r-1)-1$, hence depth $>5(r-1)-4=D_7$ at distance $<3$ from $B^k(\Lsev)$. (b) Split $e^{-\delta r}=e^{-\delta r/2}e^{-\delta r/2}$. With $l=\ln t^{-1}\ge0$: $e^{-\delta r/2}=e^{-(\delta R/2)(1+l)^\varrho}\le e^{-(\kappa+m+1)l}=t^{\kappa+m+1}$ once $(1+l)^{\varrho-1}\ge2(\kappa+m+1)/(\delta R)$, that is, for $t\le t_1$; and $b_0^n(1+l)^{n\tilde p}e^{-(\delta R/2)(1+l)^\varrho}$ is bounded in $l$; absorb this bound and $\lambda^{-m}$ by $t\le t_0\le t_1$. The second clause: $(1+l)^{n\tilde p}e^{-\beta l}$ is bounded in $l$. (c) A minimal spanning tree of the cube centres has length $\le2d(\{\square\})+4n$, so $e^{-\delta d}\le e^{2\delta n}\sum_\tau\prod_{uv\in\tau}e^{-(\delta/2)|c_u-c_v|}$ over the $n^{n-2}$ labelled trees; sum the positions leaf by leaf with $\sum_ze^{-(\delta/2)|z|}\le(1+4/\delta)^3$. The second clause is $\sum_{s\in S}\sum_\square e^{-\delta\dist(\square,s)}\le|S|(1+2/\delta)^3$.
\end{proof}

\begin{lemma}[the difference kernel as a two-variable external field]\label{lem:a07:K3}
Let $\Omega_2\subset\Omega$ be unions of large blocks, $A$ regular on $\Omega$, $\dG\defeq G_k(\Omega,A)-G_k(\Omega_2,A)$ on $\Omega_2\times\Omega_2$, and $d(x)\defeq\dist(x,\Omega_2^c)$. Let $D\ge R_0+2$, and let every evaluation point ($x$, $x'$ and the H\"older partners) have $d\ge D-1$. Then for $\sigma,\sigma'\in\{0,D,HD\}$
\[
\big|\partial^\sigma_x\partial^{\sigma'}_{x'}\dG(x,x')\big|\ \le\ C_{\dG}\,\exp\big[-\delta_{\dG}\big(d(x)+d(x')+|x-x'|\big)\big],
\]
and consequently, for localisations $h,h'$ supported at depth $\ge D$,
\[
\norm{h\,\dG\,h'}_{1,\alpha}\ \le\ C_hC_{\dG}\exp\big[-\delta_{\dG}\big(\dist(\supp h,\Omega_2^c)+\dist(\supp h',\Omega_2^c)+\dist(\supp h,\supp h')\big)\big].
\]
$C_{\dG}$ and $\delta_{\dG}$ depend on the constants of \textup{(P1)--(P3)}, on $L$ and on $\alpha$, and are free of $k$, $\eta$, $\Omega$, $\Omega_2$ and $A$. This is the two-variable, two-sided form of the domain-difference clauses \cite[(2.26)]{Ba1}, \cite[(1.20)]{Ba4}, \cite[(2.5)]{Ba3}.
\end{lemma}

\begin{proof}
\emph{Slices.} For $1\le j\le k+1$ put $G^{[k]}_j(\Omega,A)\defeq(-\Delta^N_{A,\Omega}+m^2t^2+a_jl_j^{-2}P_j(A))^{-1}$, $l_j\defeq5^{j-k}$, $P_j(A)$ the projection onto the transported constants of the scale-$j$ blocks; $G^{[k]}_k=G_k(\Omega,A)$. Rescaling $x=l_j\hat x$, $\hat A=l_j^{1/2}A$ keeps the transporters and gives $G^{[k]}_j(x,x')=l_j^{-1}\hat G_j(\hat x,\hat x')$ with $\hat G_j$ the scale-$j$ propagator in scale-$j$ units, $\hat A$ regular at scale $j$. The block identity at scale $j$ (the $\phi\phi$ block of the Schur inverse of the block operator, with Ba{\l}aban's recursion $a_{j+1}=aa_j/(aL^{-2}+a_j)$ \cite[p.~609]{Ba1}) gives, for $1\le j\le k$,
\[
S_j\defeq G^{[k]}_{j+1}-G^{[k]}_j,\qquad S_j(x,x')=a_j^2l_j^{-1}\sum_{\hat y,\hat y'\in\Omega^{(j)}}\mathsf L_j(\hat x,\hat y)\,C^{(j)}(\hat y,\hat y')\,\mathsf L_j(\hat x',\hat y')^{T},
\]
with the leg $\mathsf L_j(\hat x,\hat y)\defeq(\hat G_j\,Q_j^*(\hat A)\delta_{\hat y})(\hat x)$ and the block covariance $C^{(j)}=(\Delta^{(j)}(\Omega,\hat A)+aL^{-2}P_L)^{-1}$ of (P2). Telescoping, $G_k(\Omega,A)=G^{[k]}_1+\sum_{j=1}^{k-1}S_j$; this is Ba{\l}aban's \cite[(2.6), p.~424]{Ba3} in kernel form.

\emph{Slice bounds.} At points at distance $\ge R_0$ from $\Omega^c$, (P1) at scale $j$ applied to the bounded source $Q_j^*\delta_{\hat y}$ (sup $1$, support $B_j(\hat y)$) gives $|\partial^\sigma\mathsf L_j|\le c_0e^{-\delta_0\dist_j(\hat x,B_j(\hat y))}$ for each $\sigma$ (the depth hypothesis holds since $\dist_j\ge\dist\ge R_0$; the H\"older clause (2.24) has no restriction on the separation), and (P2) gives $|C^{(j)}(\hat y,\hat y')|\le c_2e^{-\theta_0|\hat y-\hat y'|}$. Since $|\hat x-\hat x'|\le\dist(\hat x,B(\hat y))+|\hat y-\hat y'|+\dist(\hat x',B(\hat y'))+2$, keeping half of each exponent for the sums over $\hat y,\hat y'$ (\cref{lem:a07:G}(c)) and rescaling (each $D$ is $l_j^{-1}\hat D$, each H\"older quotient $l_j^{-1-\alpha}$ times the scale-$j$ one) gives
\[
|\partial^\sigma_x\partial^{\sigma'}_{x'}S_j(x,x')|\ \le\ C_Kl_j^{-1-|\sigma|-|\sigma'|}e^{-\delta_K|x-x'|/l_j},\qquad\delta_K\defeq\tfrac12\min(\delta_0,\theta_0),
\]
with $C_K$ free of $k$, $j$, $\eta$, $\Omega$, $A$; and $|\partial^\sigma\partial^{\sigma'}G^{[k]}_1|\le C_1(L)l_1^{-1-|\sigma|-|\sigma'|}e^{-\theta_0|x-x'|/(2l_1)}$ by (P3) in scale-$1$ units (spacing $1/L$; point and bond masses of norm $L^{3/2}$; H\"older quotients over separations $\ge1/L$ bounded by $2L^\alpha\sup|D\cdot|$).

\emph{Difference of slices.} Apply the slice identity on $\Omega$ and on $\Omega_2$ with the same $A$: $\dG=\delta G^{[k]}_1+\sum_{j=1}^{k-1}\delta S_j$. By multilinearity each $\delta S_j$ is $a_j^2l_j^{-1}$ times the sum of three terms, $\delta\mathsf L\,C\,\mathsf L^T$, $\mathsf L_2\,\delta C\,\mathsf L^T$, $\mathsf L_2C_2\,\delta\mathsf L^T$, with $\mathsf L=\mathsf L_j(\Omega)$, $\mathsf L_2=\mathsf L_j(\Omega_2)$ (extended by $0$ for $\hat y\notin\Omega^{(j)}_2$), $\delta\mathsf L=\mathsf L-\mathsf L_2$, $C=C^{(j)}(\Omega)$, $C_2=C^{(j)}(\Omega_2)$ (extended by $0$), $\delta C=C-C_2$. Write $d_j\defeq l_j^{-1}d$. \emph{(i) $\delta\mathsf L$.} For $\hat y\in\Omega^{(j)}_2$, (P1)'s clause (2.26) at scale $j$ gives $c_0e^{-\delta_0(\dist(\hat x,B(\hat y))+\dist(B(\hat y),\Omega_2^c)+d_j(\hat x))}$; for $\hat y\notin\Omega^{(j)}_2$, $\delta\mathsf L=\mathsf L(\hat x,\hat y)\le c_0e^{-\delta_0\dist(\hat x,B(\hat y))}$ by (2.25), and $\dist(\hat x,B(\hat y))\ge d_j(\hat x)$ since $B(\hat y)\subset\Omega_2^c$; in all cases $|\partial^\sigma\delta\mathsf L|\le c_0e^{-(\delta_0/2)(\dist(\hat x,B(\hat y))+d_j(\hat x))}$. \emph{(ii) $\delta C$.} For $\hat y,\hat y'\in\Omega^{(j)}_2$ split $\delta C=[C_{\Omega^{(j)}}(\Omega)-C_{\Omega^{(j)}_2}(\Omega)]+[C_{\Omega^{(j)}_2}(\Omega)-C_{\Omega^{(j)}_2}(\Omega_2)]$, both restricted to $\Omega^{(j)}_2$. The first bracket is $-(PXP)^{-1}PXQX^{-1}P$ with $X=\Delta^{(j)}(\Omega)+aL^{-2}P_L$, $P=\mathbf 1_{\Omega^{(j)}_2}$, $Q=1-P$ (Schur identity); $P_L$ has no entries across $\Omega^{(j)}_2$, a union of $L$-blocks, so $PXQ=P\Delta^{(j)}(\Omega)Q$, whose kernel decays exponentially by (P2) (2.27), and both inverses decay at rate $\theta_0$; hence the bracket is $\le Ce^{-\delta_C(|\hat y-\hat y'|+\tilde d(\hat y)+\tilde d(\hat y'))}$ with $\tilde d(\hat y)\defeq\dist(\hat y,(\Omega^{(j)}_2)^c)$, because the path through $Q$ has length $\ge\tilde d(\hat y)+\tilde d(\hat y')$. The second bracket is the domain difference of the restricted block covariance, which has the same form by (P2), \cite[(1.20)]{Ba4}. If $\hat y\in\Omega^{(j)}_2$ and $\hat y'\notin\Omega^{(j)}_2$ (the only other pattern meeting a nonzero $\mathsf L_2$), $\delta C=C(\hat y,\hat y')\le c_2e^{-(\theta_0/2)(|\hat y-\hat y'|+\tilde d(\hat y))}$. \emph{(iii) Assembly.} Proceed as for $S_j$, carrying the decay factors and using $d_j(\hat x')\le d_j(\hat x)+|\hat x-\hat x'|$ and its symmetric form: each term is $\le Ce^{-\delta'(d_j(\hat x)+d_j(\hat x')+|\hat x-\hat x'|)}$, and after rescaling $|\partial^\sigma\partial^{\sigma'}\delta S_j|\le Cl_j^{-1-|\sigma|-|\sigma'|}e^{-\delta'\mathsf S/l_j}$ with $\mathsf S\defeq d(x)+d(x')+|x-x'|$. \emph{(iv) $\delta G^{[k]}_1$.} In scale-$1$ units use the layer identity $H_\Omega=(H_{\Omega_2}\oplus H_{\Omega\setminus\Omega_2})+D_\Sigma^*D_\Sigma$, $\Sigma$ the bonds of $\Omega$ with exactly one end in $\Omega_2$ and $D_\Sigma$ the covariant difference on them (valid because $P_j$ and $m^2$ are block-diagonal and both pieces are unions of blocks), which gives $(G(\Omega)f)|_{\Omega_2}-G(\Omega_2)f=-(G(\Omega)D_\Sigma^*D_\Sigma(G(\Omega_2)f\oplus0))|_{\Omega_2}$; (P3) on $\Omega$ and on $\Omega_2$ bounds the two factors pointwise at the fixed spacing, with no restriction near $\partial\Omega_2$, for each of the four derivative patterns, and the layer weight $L^{-2}$ per bond of $\Sigma$ summed gives $|\partial^\sigma\partial^{\sigma'}\delta\hat G_1|\le C(L)e^{-(\theta_0/2)(d_1(\hat x)+d_1(\hat x')+|\hat x-\hat x'|)}$; rescale with $l_1$. \emph{(v) Sum over $j$.} With $l_j=5^{-n}$, $n=k-j\ge1$, and $\mathsf S\ge2D-2\ge2$: $\sum_n5^{n(3+2\alpha)}e^{-\delta'5^n\mathsf S}\le e^{-\delta'\mathsf S/2}\sum_n5^{n(3+2\alpha)}e^{-\delta'5^n/2}\le Ce^{-\delta'\mathsf S/2}$, since $5^n\mathsf S\ge(5^n+\mathsf S)/2+5^n/2$ for $\mathsf S\ge2$. The localised form follows from the definition of the norm.
\end{proof}

\paragraph{The interaction terms.} The interaction terms of the stage of \cite[p.~579]{Ba2}, after the removals made there (the $R$-vertices, the terms of order $>\bar n$, the terms containing $\tilde B^{(k)}$, the terms localised in $B^k((\Lsev)^c)$), are a finite sum, over label multisets $\cL$ whose boxes lie within $O(1)$ of $B^k(\Lsev)$ (the terms localised in $B^k((\Lsev)^c)$ having been estimated separately on \cite[p.~579]{Ba2}, and \cite[Prop.~2.6]{Ba2} being stated for the terms localised in $B^k(\Lsev)$), over types $T$ and over orders $o\le\bar n$, of per-label full classes $\cC$ (\cref{def:a07:class}): each $\cC$ is the set of all connected graphs of that type, order and label multiset built from the vertices \cite[(1.6)--(1.15)]{Ba3} of the stage, with its Wick and Taylor weights, with the counterterm vertices \cite[(1.7)]{Ba3} carrying their fixed coefficients \cite[(1.23), (1.29)]{Ba3} and labels, \emph{including the whole-graph counterterm members} (Ba{\l}aban's $G'$ with $G_0=G$, \cite[p.~420]{Ba3}: the lone (1.7) vertex with two external legs is a member of the two-scalar-leg classes). This is the connected expansion of the effective action at fixed order and labels; the one-particle-irreducibility of \cite[p.~430]{Ba3} concerns the counterterm-defining classes, whose coefficients are fixed, not the members. The terms may also be finite products of such classes contracted by fixed many-variable external fields \cite[p.~419]{Ba3}. The scalar lines are $G_A\defeq G_k(\OmA,\Bt)$; the external scalar fields are $E_A\defeq a_kG_Af$ with $\norm f_\infty\le Cp\lambda(t)^{-1/4}$, and the other external fields have norms $\le Cp^a\lambda(t)^{-b}$ with $a,b$ bounded by $\bar n$. All vertices lie in $B^k(\Lsev)\subset\OmC$.


% a07_k01steps_c: last part of app:a07_k01steps (input from sec/a07_k01steps.tex): the regrouping into full classes of the cut components, their coverage by Balaban's renormalisation bound, and the theorem.

\begin{lemma}[the expansion]\label{lem:a07:M}
Put $G_C\defeq G_k(\OmC,\Bt)$ ($=G_k(\OmC,B^{(k+1),\eta})$ by \cref{lem:a07:C51}), $\dG\defeq G_A-G_C$ on $\OmC\times\OmC$ with $G_C$ extended by $0$, $E_C\defeq a_kG_Cf$ and $\dE\defeq E_A-E_C=a_k\dG f_{\mathrm{in}}+a_kG_Af_{\mathrm{out}}$ with $f_{\mathrm{in}}\defeq\mathbf 1_{\OmC}f$, $f_{\mathrm{out}}\defeq f-f_{\mathrm{in}}$. For a class $\cC$, or a product of classes contracted by fixed fields,
\[
\cC[G_A;E_A]-\cC[G_C;E_C]\ =\ \sum_{(n,m)\ne(0,0)}\cC^{(n,m)},
\]
where $\cC^{(n,m)}$, the coefficient of $s^nu^m$ in $\cC[G_C+s\,\dG;E_C+u\,\dE]$, is the sum over $\Gamma\in\cC$, over sets $\Lambda$ of $n$ scalar lines of $\Gamma$ and over sets $M$ of $m$ $E$-legs of $\Gamma$, of the value of $\Gamma$ with $\dG$ on $\Lambda$, $G_C$ on its other scalar lines, $\dE$ on $M$ and $E_C$ on its other $E$-legs.
\end{lemma}
\begin{proof}
A class value is a polynomial, multilinear in the scalar line kernels and in the $E$-legs (Wick expansion); put $s=u=1$. The counterterm coefficients are vertex data built from torus propagators at zero background \cite[(1.29)--(1.31)]{Ba3} and are not expanded: the replacement of propagators acts on lines only.
\end{proof}


\begin{definition}[cut skeletons and $\sigma$-classes]\label{def:a07:sigma}
Fix $\cC$ (type $T$, order $o$, labels $\cL$) and $(n,m)$. A \emph{cut skeleton} $\sigma$ consists of a partition $\cL=\cL_1+\dots+\cL_\rho$ into non-empty sub-multisets with orders $o_1+\dots+o_\rho=o$, $o_c\ge0$; a connected multigraph $Q$ on $\{1,\dots,\rho\}$ with $n$ edges (loops and multiple edges allowed); and an assignment of the external legs of $\cC$ to the components, each $E$-leg marked $\dE$ or $E_C$ with $m$ marked $\dE$. For each component $c$, $\cC_{\sigma,c}$ is the per-label full class of all connected configurations with labels $\cL_c$ and order $o_c$, whose external legs are the legs assigned to $c$ together with $\deg_Q(c)$ scalar \emph{stub} legs, with scalar lines $G_C$ and vector lines as in $\cC$, \emph{members without internal line and of order $0$ included}: equivalently, the coefficient of the monomial of the labels $\cL_c$ in the stub-Taylor expansion of the full connected generating function $\ln\int e^{-V(\cdot+\chi+\psi)}d\mu_{G_C}$, $\psi$ the stub field. The stub and $E$-legs may sit at any scalar slot of any vertex, differenced or not; the class sums over all of them. The \emph{$\sigma$-class} is the contraction $\Sigma_\sigma(\cC)\defeq\big\langle\bigotimes_c\cC_{\sigma,c},\bigotimes_{e\in Q}\dG_e\big\rangle$: an edge $\{c,c'\}$ contracts one stub of $c$ with one stub of $c'$ by $\dG$, a loop two stubs of $c$. It carries the Taylor weight $w_\sigma$, with $0\le w_\sigma\le C(\bar n)$.
\end{definition}

\begin{lemma}[regrouping into $\sigma$-classes]\label{lem:a07:R}
$\cC^{(n,m)}=\sum_\sigma w_\sigma E(\Sigma_\sigma(\cC))$, the sum over the cut skeletons of $(\cC;n,m)$.
\end{lemma}
\begin{proof}
\emph{The bijection.} To $(\Gamma,\Lambda,M)$ assign the components of $\Gamma\setminus\Lambda$ (connected through the kept lines, vector lines included, since $\Lambda$ consists of scalar lines), their label multisets and orders, the multigraph $Q$ that $\Lambda$ induces on them (connected, since $\Gamma$ is), the assignment of the external legs, and for each $c$ the restriction $\Gamma_c$ of $\Gamma$ to $c$ with the ends of $\Lambda$-lines at $c$ turned into stubs. Then $\Gamma_c\in\cC_{\sigma,c}$: it is connected, with labels $\cL_c$, order $o_c$ and its external legs; it may be a single vertex all of whose legs are stubs or external fields (when every line at that vertex is in $\Lambda$), or a path of order-$0$ vertices (when $\Lambda$ cuts out a stretch of an expanded line), which is why these members are in $\cC_{\sigma,c}$ (with \cite[p.~415]{Ba3}'s ``there is at least one internal line'' the map would not land in the product). Conversely, given $\sigma$, members $\Gamma_c\in\cC_{\sigma,c}$ and a pairing of all stubs compatible with $Q$, joining the paired stubs by lines gives a connected graph of type $T$, order $o$ and labels $\cL$, hence a member of the full class $\cC$, with $\Lambda$ the set of joining lines; the two maps are inverse. \emph{The weights.} In the labelled form with vertex labels $1,\dots,N$ and weight $\frac1{N!}\prod_v(-c_v)$, the components take disjoint sets of vertex labels, in $N!/\prod_cN_c!$ ways, so the weight is $\prod_c\frac1{N_c!}\prod_{v\in c}(-c_v)$ times the symmetry factor of identical components, parallel edges and loops of $Q$, which is that of the Taylor coefficients. \emph{Generating-function form} (the same identity, weights included): with $W(K;\chi)$ the connected generating function at fixed labels, scalar covariance $K$ and external field $\chi$, $W(G_C+s\,\dG;\chi)=\ln\E_{s\dG}\big[\exp W(G_C;\chi+\psi)\big]$, a formal Gaussian convolution in the stub field $\psi$ (the moments of covariance $K_1+K_2$ are the convolution of those of $K_1$ and $K_2$, an identity of polynomials in the kernels; no positivity of $\dG$ is used). By the linked-cluster theorem \cite{GJ} the right side is a sum of connected graphs whose vertices are the $\psi$-Taylor coefficients of $W(G_C;\chi+\psi)$, the classes $\cC_{\sigma,c}$ with their stubs, and whose lines are $\dG$; taking the coefficient of the labels and of $s^n$, with $\chi=E_C+u\,\dE$, and then of $u^m$, gives the claim. \emph{Products.} For a product of classes contracted by fixed many-variable fields \cite[p.~419]{Ba3}, apply the above factor by factor; the fixed fields stay external fields of the components they touch and join the components in step 5 below exactly as the $\dG$ fields do.
\end{proof}

The classes $\cC_{\sigma,c}$ are full per-label classes, not orbits: the finer groupings (\cref{lem:a07:orb}) fix slot data, which is needed where a bound depends on slot types; here only the $\norm\cdot_{1,\alpha_0}$ norms of $\dG$ and $\dE$ are consumed, which control both slot types (\cref{lem:a07:K3} has $\sigma,\sigma'\in\{0,D,HD\}$ in each variable), so the full grouping is admissible, and it is the form in which the renormalisation bound is available at a nonzero background. \Cref{lem:a07:R} is Ba{\l}aban's own step at \cite[p.~432]{Ba3} for the lines $\delta G_k$ (``we replace these lines by pairs of external legs \dots\ $G'$ may be unconnected \dots\ it can be represented as a sum of classes $G'_{\mathrm{ren}}$''), which for full classes is exact.

\begin{lemma}[the $\sigma$-classes at a nonzero background]\label{lem:a07:S}
For every cut skeleton $\sigma$ and component $c$:
\begin{enumerate}[label=\textup{(\alph*)},nosep,leftmargin=*]
\item $\cC_{\sigma,c}$ is a family of the form of \cite[Prop.~1]{Ba3} on the domain $\OmC$ at the background $\Bt$: vertices \cite[(1.6)--(1.15)]{Ba3} localised in $B^k(\Lsev)$, scalar lines $G_k(\OmC,\Bt)$, vector lines as printed, and external fields of several variables (the stub kernels $\dG$, ``treated as an external field'' \cite[pp.~414, 419]{Ba3}; the fields $\dE$, $E_C$; the remaining fields of $\cC$).
\item Ba{\l}aban's classes are sets of graphs built from the vertices of \cite[(1.5)]{Ba3}, and $\Bt$ enters only the line operators and the transporters of \cite[(1.12)--(1.15)]{Ba3}; so $G_{\mathrm{ren}}(G)$ is the same set of graphs at $\Bt=0$ and at $\Bt\ne0$. Every completion keeps type, order and label multiset; hence $G_{\mathrm{ren}}(G)\subset\cC_{\sigma,c}$ for every member $G$, and $\cC_{\sigma,c}$ is the union of the $G_{\mathrm{ren}}(G)$ over its members.
\item $\cC_{\sigma,c}$ is closed under the exchange of the two external scalar legs of every two-scalar-leg or $\phi\phi A$ subgraph, the whole graph included (\cref{lem:a07:full}(c)), exchanges that move a stub or change its slot type included; and its Taylor coefficients in the fluctuation of the background are again per-label full classes, every touched two-scalar-leg subgraph having its mirror.
\item Split $\cC_{\sigma,c}=\cC^{\mathrm{ren}}_{\sigma,c}+\cC^{0}_{\sigma,c}$, where $\cC^{\mathrm{ren}}_{\sigma,c}$ consists of the members of order $\ge1$ which are graphs in the sense of \cite[p.~415]{Ba3} together with the lone counterterm vertex \cite[(1.7)]{Ba3} (the counterterm member of the whole-graph pairing), and $\cC^0_{\sigma,c}$ of the bare single vertices with every leg a stub or an external field and of the order-$0$ paths of $n=0$ vertices joined by $G_C$ lines. Then $\cC^{\mathrm{ren}}_{\sigma,c}$ is a per-label full class of the form of \cite[Prop.~1]{Ba3}, to which the bound \cite[(1.33)]{Ba3} applies at block factor $5$ with constants ``independent of $\eps$, $k$, the domains $\Omega,\Omega_1,\Omega_2$, the vector field $\tilde B$'' \cite[p.~421]{Ba3} (\cref{thm:a06:BALA}; the unprinted steps of that proposition are those listed in Step~3 of the proof of \cref{thm:a06:T1}), with the localisation $h$ of each external field equal to $1$ on the label boxes of every vertex at which the field sits in some member of the class, the class being split by the cluster at which each leg sits (at most $V(\bar n)^{\#\mathrm{legs}}$ pieces). Each member of $\cC^0_{\sigma,c}$ obeys a bound of the form of \cite[(1.33)]{Ba3} directly: a bare vertex is at most its coupling times $\prod_i\norm{h_i\Phi_i}_{1,\alpha_0}$ summed over its cube or block (\cite[p.~421]{Ba3}: ``the vertex itself with all external legs \dots\ is treated separately''); an order-$0$ path is a tree whose only pieces of degree $\le0$ are the pairs $VG_CV$ of \cite[(2.4)]{Ba3}, so that \cite[Prop.~2.1]{Ba3} applies to it: each pair is treated by \cite[(2.8)--(2.9)]{Ba3} at one vertex, the derivative moved onto a stub, an $E$-leg or the external vector leg (allowed since the norm \cite[(1.32)]{Ba3} contains one derivative, and \cref{lem:a07:K3} covers the stub), after which all degrees are positive. No member of $\cC^0_{\sigma,c}$ is a partner that the bound needs: partners have the same order $\ge1$, and a zero-line member of order $\ge1$ other than a counterterm vertex is a (1.6) vertex with four external legs, alone in its class. The factors of the background in these bounds are $O(e(t)p)$ \cite[p.~412]{Ba3}.
\end{enumerate}
\end{lemma}
\begin{proof}
(a) By construction and \cref{lem:a07:C51}. (b) At \cite[p.~414]{Ba3} ``the vertices of (1.5) are given by (1.6)--(1.8), (1.10), (1.13), and (1.14) with $n'=0$ and $\tilde B=A^{(k)}$'', and the lines are ``$G_k(\Omega,\tilde B)$, $G_k(\Omega_2,\tilde B)$, $\delta G_k(\Omega,\Omega_2,\tilde B)$ or the operator (1.16)''; the completions (the counterterm graph of a divergent two-scalar-leg subgraph \cite[p.~420]{Ba3}, the permutation \cite[p.~429]{Ba3}, the two-vector-leg class in its maximal form \cite[p.~431]{Ba3}) keep type, order and labels (``localizations in these graphs should be in agreement''), and a per-label full class contains all graphs of its type, order and labels. (c) The first clause is \cref{lem:a07:full}(c); the second is the identity of generating functions in its proof: the expansion in the fluctuation $\tilde B'$ of the background is the Taylor expansion of the full connected generating function, whose coefficients are full classes with $\tilde B'$-legs carrying their partition labels, and the mirror of a touched two-scalar-leg subgraph is an exchange of the first clause. (d) The members of $\cC^{\mathrm{ren}}_{\sigma,c}$ form the per-label full class of \cite[Prop.~1]{Ba3}'s form by (a)--(c); the lone (1.7) vertex is the counterterm member of the whole-graph pairing (\cite[p.~417]{Ba3}: ``the same graph $G$ but with both external legs localized in $x$''). The localisation rule for a field that sits at several vertices of a class is the reading of \cite[p.~420]{Ba3} (``$h=1$ on $\mathrm{box}(v)$'') for classes; per piece the stub bound of step 4 below holds with the piece's $h$, whose supports are at least as far apart as the original ones. The bounds on $\cC^0_{\sigma,c}$ are as stated: a bare vertex carries no line and no loop; an order-$0$ path consists of $n=0$ vertices, each with at most two scalar legs, so it is a tree, and \cite[(2.8)]{Ba3} is the one-derivative move of \cite[Sect.~2]{Ba3} at a vertex of a (2.4)-type pair. The last two clauses are read off the orders.
\end{proof}

\begin{theorem}[\BALint: the interaction-term half of {\cite[Prop.~2.6]{Ba2}} at every background]\label{thm:a07:dom}
In the setting above, for every $\kappa>0$ there are $C_\kappa<\infty$ and $t_0(\kappa)>0$ such that for $t\le t_0(\kappa)$
\[
\sum_{\cL}\sum_{\cC}\sum_{(n,m)\ne(0,0)}\sum_{\sigma}w_\sigma\big|E\big(\Sigma^{(n,m)}_\sigma(\cC)\big)\big|\ \le\ C_\kappa\,t^{\kappa}\,|\Lsev|,
\]
the sums over every localisation $\cL$ of the stage (vertex cubes in $B^k(\Lsev)$, label boxes within $O(1)$ of it), every class $\cC$ with labels $\cL$, every $(n,m)\ne(0,0)$ and every cut skeleton $\sigma$. By \cref{lem:a07:M,lem:a07:R} the left side dominates $\sum_\cL\sum_\cC|\cC[G_A;E_A]-\cC[G_C;E_C]|$: this is \cite[Prop.~2.6]{Ba2} for the interaction terms, with error $O((L^k\eps)^\kappa)|\Lsev|$ for every $\kappa$, at every background $\Bt$, $\Bt=0$ included; $C_\kappa$ and $t_0(\kappa)$ depend on $\kappa$, on $(R,\varrho,b_0,\tilde p,\bar n,\alpha_0)$, on the couplings and on the constants of \textup{(P1)--(P3)} and of \cite[(1.33)]{Ba3} only. This proves \cref{hyp:k01:BALint}.
\end{theorem}

\begin{proof}
\emph{Step 1 (domain change only).} By \cref{lem:a07:C51} the replacement is $G_A\to G_C$ at the one background $\Bt$; $\OmC\subset B^k(\Lambda^{(k)}_0)\subset\OmA$, and $\Bt$ is regular on $\OmA$.

\emph{Step 2 (expansion and regrouping).} By \cref{lem:a07:M,lem:a07:R} the difference is $\sum_{(n,m)\ne(0,0)}\sum_\sigma w_\sigma E(\Sigma^{(n,m)}_\sigma(\cC))$.

\emph{Step 3 (depth).} Every vertex, every label box and every localisation support lies within $2+\frac14+\eta<3$ of $B^k(\Lsev)$, hence at depth $\ge D_7=5(r-1)-4$ in $\OmC$ by \cref{lem:a07:G}(a), with $D_7\ge R_0+2$ by \eqref{eq:a07:D7}. The domain hypothesis of \cite[Prop.~1]{Ba3} for $\Omega_2=\OmC$ (``the same conditions as $\Omega$'' \cite[p.~414]{Ba3}, including $\dist(\supp\tilde A,\partial\Omega_2)>2r$ \cite[p.~412]{Ba3}) holds for the $\tilde A$-legs actually present: with $\tilde A=\theta_{k+1}A'^{(k)}$ the raw support reaches $\partial\OmC$, but after \cite[p.~579]{Ba2} every vertex and every localised leg lies within $O(1)$ of $B^k(\Lsev)$, at depth $\ge D_7>2r$ (since $r\ge R\ge R_0+3>3$); multiplying $\tilde A$ by a cut-off equal to $1$ on that neighbourhood and supported at depth $>2r$ changes no term.

\emph{Step 4 (external fields).} For the localisations $h,h'$ of the two ends of an edge of $Q$ ($h=1$ on the label boxes of the vertices at which the stub can sit in its component, supported in their $\frac14$-neighbourhood), \cref{lem:a07:K3} with $(\Omega,\Omega_2,A)=(\OmA,\OmC,\Bt)$ gives
\[
\norm{h\,\dG\,h'}_{1,\alpha_0}\ \le\ C_hC_{\dG}\,e^{-2\delta_{\dG}D_7}\,e^{-\delta_{\dG}\dist(\supp h,\supp h')},
\]
with the mixed clauses $\sigma,\sigma'\in\{0,D,HD\}$ that the joining of components needs. For $\dE$: $\norm{h\,\dE}_{1,\alpha_0}\le2c_0e^{-\delta_0D_7}\norm f_\infty$, since (P1) (2.26) bounds $a_k\dG f_{\mathrm{in}}$ and (2.25) bounds $a_kG_Af_{\mathrm{out}}$ at points at distance $\ge D_7$ from $\supp f_{\mathrm{out}}\subset\OmC^c$. The norms of $E_C$ and of the other external fields are $\le Cp^a\lambda(t)^{-b}$ ((P1) (2.24)--(2.25) on $\OmC$ at depth $D_7$).

\emph{Step 5 (the renormalisation bound per $\sigma$-class).} By \cref{lem:a07:S}(d), \cite[(1.33)]{Ba3} or the direct bound applies to each $\cC_{\sigma,c}$ at fixed labels with the external fields of Step 4. The components are joined along a spanning tree of $Q$ by the decay of the stub kernels, the non-tree edges and loops entering as many-variable external fields (\cite[(3.2), p.~432]{Ba3}: ``the exponential factor in (1.33) is obtained from the estimates of external fields in (3.2)''). With $e_k,\lambda_k\le1$,
\[
\big|E\big(\Sigma^{(n,m)}_\sigma(\cC);\cL\big)\big|\ \le\ O(1)\,e^{-\delta''d(\cL)}\,\big(C'e^{-2\delta_{\dG}D_7}\big)^{n}\big(2c_0Cp\lambda^{-1/4}e^{-\delta_0D_7}\big)^{m}\big(Cp^a\lambda^{-b}\big)^{n_{\mathrm{ext}}},
\]
where $d(\cL)$ is the tree length of all localisation cubes over all components and $\delta''\defeq\frac12\min(\delta_0,\delta_{\dG},\delta_{\mathrm{ren}})$, $\delta_{\mathrm{ren}}$ the rate of \cite[(1.33)]{Ba3} (called $\delta_0$ there; it is not the $\delta_0$ of \textup{(P1)}): the trees inside the components come from (1.33), the edges between components from the stub decay, and the $\frac14$-neighbourhoods change the tree length by at most a constant times the number of cubes, a constant of $\bar n$ absorbed in $O(1)$.

\emph{Step 6 (size).} Since $n+m\ge1$ the product carries $e^{-\delta_*D_7}$ at least once, $\delta_*\defeq\min(2\delta_{\dG},\delta_0)$, and $D_7\ge r$; by \cref{lem:a07:G}(b), $p^a\lambda^{-b}e^{-\delta_*D_7}\le t^\kappa$ for $t\le t_0(\kappa)$, the background factors $O(e(t)p)$ of \cref{lem:a07:S}(d) being absorbed the same way.

\emph{Step 7 (sum).} Fix one cube in $B^k(\Lsev)$ ($|\Lsev|$ unit cubes, and every cube lies there); \cref{lem:a07:G}(c) sums the others with $e^{-\delta''d}$. The numbers of types, of orders $\le\bar n$, of pairs $(n,m)$ and of skeletons $\sigma$ are bounded by constants of $\bar n$ ($\rho\le V(\bar n)$, $n$ at most the number of scalar lines, $Q$ and the leg assignments finite).
\end{proof}

\begin{remark}\label{rem:a07:dom}
(i) The proof uses no telescoping in a fixed line order and no forest families: the terms with one difference line, grouped per line, carry their counterterm partners but not the two-vector-leg and permutation partners of the cut graph (\cref{rem:a07:orb}(iv)), and a lone one-loop two-vector-leg subgraph is $O(\eta^{-1})$ per graph, which $e^{-2\delta_{\dG}D_7}=O(t^\kappa)$ does not cure; the regrouping of \cref{lem:a07:R} restores the complete classes, at every background. (ii) The corresponding replacement in the polynomial $Q^{(k)}$ of \cite[p.~579]{Ba2}, the other half of \cite[Prop.~2.6]{Ba2}, is not used through \cref{hyp:k01:BALint}. (iii) On the torus skew in the rotated frame (the instance consumed by the rotation part through \cite[Thm.~3.1]{K1} there), \cref{lem:a07:M,lem:a07:R,lem:a07:S} are combinatorics without geometry, \cite[(1.33)]{Ba3} transfers as the other torus estimates do, and the proof of \cref{lem:a07:K3} uses no cubic symmetry, its inputs (P1)--(P3) being available on that torus: \cref{lem:a07:Ptorus,lem:a07:K3torus} below carry (P1)--(P3), \cref{lem:a07:K3} and \cref{thm:a07:dom} to every lattice torus fitted to the blocks, these tori included.
\end{remark}

% k0x_fill_skew: the inputs (P1)--(P3), the difference-kernel bound of lem:a07:K3 and \BALint\ on a lattice torus fitted to the blocks
% (input from sec/a07_k01steps_c.tex after rem:a07:dom). This is the transfer consumed by the rotation part on the tori of rot:hyp:F.

\subsection{The difference-kernel bound on a lattice torus fitted to the blocks}\label{ssec:a07:torus}

The rotation part applies King's \cite[Thm.~3.1]{K1}, and through it \BALint, on the tori $T_1=\DJ$ and $T_2=R^{-1}\DJ$ of \cref{rot:hyp:F}, each carrying both lattices $\Lk$ and $R^{-1}\Lk$ of \cref{rot:def:tori}(b). In the frame of the lattice $R^{-1}\Lk$ the torus $\DJ$ is the quotient of that lattice's own axis lattice by $R(50MJ\,\Z^3)$, and in the frame of $\Lk$ the torus $R^{-1}\DJ$ is the quotient by $R^{-1}(50MJ\,\Z^3)$: lattices of translations which are not rectangular in those frames. The statements of this appendix were written for subsets of $\eta\Z^3$. This subsection records that their proofs use the global geometry only through the two properties of \cref{def:a07:btorus}, which the tori of the rotation part have, so that \cref{lem:a07:CT,lem:a07:K3} and \cref{thm:a07:dom} hold on them with constants that do not depend on the torus.

\begin{definition}[a lattice torus fitted to the blocks]\label{def:a07:btorus}
In scale-$k$ units let $T=\eta\Z^3/\Gamma$, where $\Gamma\subset\Z^3$ is a lattice of rank $3$ (so $T$ carries the unit blocks and, since $\Gamma\subset\Z^3$, the blocks of every scale $j\le k$, of side $l_j=5^{j-k}$). Let $M_b\in\N$ be the side of a large block. $T$ is \emph{fitted to the blocks} if $\Gamma\subset M_b\Z^3$ and its sup-systole $\mathrm{sys}_\infty(T)\defeq\min\{|\gamma|_\infty:\gamma\in\Gamma\setminus\{0\}\}$ is at least $3M_b$. Distances on $T$ are the quotient sup-distances $\dist_T(x,x')\defeq\min_{\gamma\in\Gamma}|x-x'-\gamma|_\infty$; $\Omega\subset T$ denotes a union of large blocks, $T$ itself included, and ``regular on $\Omega$'' is \cite[(2.23)]{Ba1} read on the bonds of $\Omega$.
\end{definition}
If $T$ is fitted to the blocks in scale-$k$ units, so is its image in scale-$j$ units for every $j\le k$: the lattice becomes $l_j^{-1}\Gamma\subset l_j^{-1}M_b\Z^3\subset M_b\Z^3$ (as $l_j^{-1}=5^{k-j}\in\N$) and the systole grows by the factor $l_j^{-1}$. For the tori of the rotation part, in the frame of either lattice, the lattice of translations is $50MJ\,\Z^3$ or $R^{\pm1}(50MJ\,\Z^3)\subset10MJ\,\Z^3$ in physical units (\cref{rot:lem:geom}(b)). At Ba{\l}aban's step $k$ on the lattice of spacing $\eps_K=5^{-K}$ the scale-$k$ unit is $t=5^k\eps_K=5^{k-K}\le1$, so in scale-$k$ units these lattices lie in $10MJ\,5^{K-k}\,\Z^3$ and the systole is at least $10MJ\,5^{K-k}$. They are therefore fitted to the blocks at every step whenever the large-block side $M_b$ divides $10M$ and $3M_b\le10MJ$; in King's flow, the one run on these tori through \cite[Thm.~3.1]{K1}, the large block is King's and $M_b=M$.

\begin{lemma}[the printed inputs on a torus fitted to the blocks]\label{lem:a07:Ptorus}
Let $T$ be fitted to the blocks, $\Omega\subset T$ a union of large blocks and $A$ regular on $\Omega$. Then the bounds \textup{(P1)} (\cite[Prop.~2.1, (2.24)--(2.26)]{Ba1}, \cite[(1.9)--(1.12)]{Ba4}), \textup{(P2)} (\cite[Props.~2.2--2.3]{Ba1}, \cite[(1.13)--(1.20)]{Ba4}) and the positivity \cite[(1.8)]{Ba4} hold for the operators $G_k(\Omega,A)$, $\Delta^{(k)}(\Omega,A)$, $C^{(k)}_\Lambda(\Omega,A)$ on $T$, with the distances of \cref{def:a07:btorus}, with constants $c_0$, $\delta_0$, $R_0$, $c_2$, $\theta_0$, $\gamma$ that depend on $d$, $a$, $M_b$ and $L$ ($c_0$ also on $\alpha$) and not on $\Gamma$, $k$, $\Omega$, $\Lambda$ or $A$. For $\Omega=T$ the restriction to points at distance $\ge R_0$ from $\Omega^c$ is void. The same holds in scale-$j$ units for every $j\le k$.
\end{lemma}
\begin{proof}
Ba{\l}aban's proofs \cite[Sect.~2, pp.~574--586]{Ba4} are written for a union $\Omega$ of large blocks in $\eta\Z^d$ (``the torus identified with a rectangular parallelepiped''). They use the geometry of the ambient lattice in exactly the following places, each of which we verify on $T$.

\emph{(1) The covering cubes.} Let $v_j$ run through the vertices of the large-block grid of $T$ and $C_j\defeq v_j+[-M_b,M_b)^3$, the union of the $2^3$ large blocks with vertex $v_j$ (\cite[p.~578]{Ba4}: ``a cube with center $Mj$ and size equal to $2M$''). Since $2M_b<\mathrm{sys}_\infty(T)$, the projection $\eta\Z^3\to T$ is injective on each lift of $C_j$; it preserves bonds, unit blocks and the contours $\Gamma^{(k)}_{y,x}$ of the averaging operators, which lie inside blocks \cite[(1.4)]{Ba4}. Hence $\Pi_j\defeq\Omega\cap C_j$ is a union of at most $2^3$ large blocks in an embedded cube, the constant gauge $\lambda_0(x)=A_0\cdot(x-c_j)$ is defined on $C_j$, and \cite[Lemmas~2.1, 2.2, 2.4]{Ba4}, which are statements about $G_k(\Pi_j,A_j)$ on a parallelepiped of $\eta\Z^d$ with a local field $A_j$ built from $A|_{C_j}$, apply with their constants, which depend on $d$, $M_b$, $\alpha$ and the bare $a$ only.

\emph{(2) The partition of unity.} Let $h\in C^\infty_c((-\frac23,\frac23))$ with $h=1$ on $[-\frac13,\frac13]$ and $\sum_{n\in\Z}h(\cdot-n)^2=1$, and $h_j(x)\defeq\prod_\mu h((x_\mu-(v_j)_\mu)/M_b)$ \cite[p.~575]{Ba4}. Because $\Gamma\subset M_b\Z^3$, the family $\{h_j\}$ is $\Gamma$-periodic, descends to $T$, and $\sum_jh_j^2=1$ on $T$ with $\supp h_j\subset C_j$; so \cite[(2.2)--(2.9)]{Ba4} hold on $T$ as on $\eta\Z^d$ (if $\Omega=T$, then $\partial\Omega=\emptyset$ and (2.6) is trivial).

\emph{(3) The counts.} Each $j$ has at most $3^3$ neighbours and each point lies in at most $2^3$ of the supports, so the remainder of the parametrix has $\norm R_{2,2}\le c_dO(M_b^{-1})<1$ as on \cite[p.~577]{Ba4}, and the number $2^d(3^d)^{n-1}2^d$ of walks of length $n$ \cite[p.~579]{Ba4} is unchanged.

\emph{(4) The metric.} Consecutive cubes of a walk have centres at quotient distance $\le M_b$, so a walk of length $n$ from a cube meeting $\supp f$ to a cube containing $x$ has $n\ge\dist_T(x,\supp f)/M_b-2$; the geometric series \cite[(2.22)]{Ba4} then gives $\delta_0$, $c_0$ as on $\eta\Z^d$. The definition \cite[(2.19)]{Ba4} of $R_0$ and the implication ``$\dist_T(\{x,x'\},\Omega^c)\ge R_0$ implies that the first $n_0$ cubes of the walk lie in $\Omega$'' use only this. The cancellation of the interior-cube terms in $G_k(\Omega,A)-G_k(\Omega_0,A)$, which gives the two-sided factor of \cite[(2.26)]{Ba1}, uses (1) and (2). If $\Omega=T$, every $\Pi_j=C_j$ is a full cube and no restriction on $x$, $x'$ is needed \cite[p.~611]{Ba1}.

\emph{(5) Positivity.} For \cite[(1.8)]{Ba4}, $P_k(A)$ is block-diagonal, so dropping the bonds between unit blocks bounds the form from below by a sum over the unit blocks $\Delta$ of $\langle\phi,(-\Delta^N_{A,\Delta}+a_kP_k(A))\phi\rangle_\Delta$; on each block the gauge of (1) and the expansion in the field reduce the bound to $A=0$ \cite[p.~580]{Ba4}, where it is the Poincar\'e inequality on a unit block, which embeds in $T$. So $\gamma$ is Ba{\l}aban's.

The block operators of (P2) are built from $G_k(\Omega,A)$ and the averaging operators, and the proofs of \cite[Props.~2.2--2.3]{Ba1}, \cite[Sect.~3]{Ba4} use the same cubes, partition and walk and no further property of the ambient lattice; the restricted covariances $C^{(k)}_\Lambda$ for a union $\Lambda$ of large blocks are handled by the same walk on $\Lambda$. Nothing else in \cite[Sects.~2--3]{Ba4} refers to the global geometry. The last sentence of the lemma follows because the rescaled torus is fitted to the blocks (the remark after \cref{def:a07:btorus}) and the statements are the same at every scale.
\end{proof}

\begin{lemma}[\cref{lem:a07:CT,lem:a07:K3} and \BALint\ on a torus fitted to the blocks]\label{lem:a07:K3torus}
Let $T$ be fitted to the blocks. Then:
\begin{enumerate}[label=\textup{(\alph*)},nosep,leftmargin=*]
\item \cref{lem:a07:CT} holds for subsets of $T$ with the quotient metric, with the same constants.
\item \cref{lem:a07:K3} holds for unions $\Omega_2\subset\Omega$ of large blocks of $T$ and $A$ regular on $\Omega$, with $\dist_T$ in place of $\dist$ and with $C_{\dG}$, $\delta_{\dG}$ depending on $d$, $a$, $M_b$, $L$ and $\alpha$ only.
\item \cref{thm:a07:dom} holds on $T$, with the propagator domains $\OmA\supset\OmC$ and the background $\Bt$ of Ba{\l}aban's stage on $T$ (\cite[Sect.~2]{Ba2} run on $T$, as in \cref{thm:a06:BALA}) and with the renormalisation bound \cite[Prop.~1, (1.33)]{Ba3} on $T$; its constants $C_\kappa$, $t_0(\kappa)$ are those of the statement on $\eta\Z^3$ and do not depend on $\Gamma$.
\end{enumerate}
In particular all three hold on the tori $T_1$, $T_2$ of \cref{rot:hyp:F}, in the frame of either lattice.
\end{lemma}
\begin{proof}
(a) The proof of \cref{lem:a07:CT} uses the weight $\rho=\dist(\cdot,\supp f')$ only through two properties: $\rho$ is $1$-Lipschitz along bonds (the ends of a bond are at distance $\eta$), and each site is the final point of $d$ bonds. Both hold on $T$ with $\rho\defeq\dist_T(\cdot,\supp f')$, since $\dist_T(\cdot,S)$ is $1$-Lipschitz for the quotient metric and the bond structure is that of $\eta\Z^3$ locally. The lower bound $T\ge\gamma$ it needs is \cite[(1.8)]{Ba4} on $T$ (\cref{lem:a07:Ptorus}). The constants are unchanged.

(b) The proof of \cref{lem:a07:K3} consists of the slice decomposition $G_k(\Omega,A)=G^{[k]}_1+\sum_{j=1}^{k-1}S_j$, an operator identity which needs only that the blocks of every scale $j\le k$ partition $\Omega$ (true on $T$: a union of large blocks is a union of unit blocks of every scale $j\le k$, and these descend to $T$ because $\Gamma\subset M_b\Z^3\subset l_j\Z^3$); the rescaling to scale-$j$ units, under which $T$ stays fitted to the blocks; the slice bounds, which are (P1) at scale $j$ for the legs, (P2) at scale $j$ for the block covariances and (P3) at scale $1$, all available on $T$ with $\Gamma$-free constants by \cref{lem:a07:Ptorus} and (a); the sums over the block centres $\hat y,\hat y'$, which use \cref{lem:a07:G}(c) in the form $\sum_{\square}e^{-\delta\dist_T(\square,S)}\le C(\delta)|S|$, true on $T$ because the projection of the cubes of $\eta\Z^3$ onto those of $T$ is onto, so the sum on $T$ is at most the sum on $\eta\Z^3$ over a lift of $S$; and the triangle inequality for the metric. Every step is therefore available on $T$ with the quotient metric, and the proof goes through word for word. The constants $C_{\dG}$, $\delta_{\dG}$ are built from those of (P1)--(P3), $L$ and $\alpha$.

(c) In the proof of \cref{thm:a07:dom}: \cref{lem:a07:M,lem:a07:R,lem:a07:S} are statements about the expansion and regrouping of graph sums and about the sets of graphs forming a class; they do not refer to the geometry. \cref{lem:a07:C51} uses only that the three operators depend on the field through the bonds of their domain and on contours inside blocks. \cref{lem:a07:G}(a) uses the construction of the layers $\Lambda^{(k)}_i$ by sup-neighbourhoods of unions of large blocks and sup-metric geodesics, which exist in the quotient metric (a shortest path lifts to $\eta\Z^3$); (b) is a scalar inequality; (c) holds on $T$ as shown in (b) above, the number of cubes of $T$ at a given tree distance from a fixed cube being at most the number on $\eta\Z^3$. Steps 1--3 of the proof use (a) of \cref{lem:a07:G} and the definitions of the stage; Step 4 uses (b) above and (P1) on $T$; Step 5 uses \cite[Prop.~1, (1.33)]{Ba3} on $T$, which is Ba{\l}aban's renormalisation bound for a torus aligned with the blocks (\cref{thm:a06:BALA}); Steps 6 and 7 use \cref{lem:a07:G}(b), (c). The constants are those of the statement on $\eta\Z^3$.

For the tori of the rotation part, both lattices of translations are contained in $10MJ\,\Z^3$ in the frame in question (\cref{rot:lem:geom}(b)), so the tori are fitted to the blocks with $M_b=M$ by the remark after \cref{def:a07:btorus}.
\end{proof}

\begin{remark}\label{rem:a07:torus}
(i) The two conditions of \cref{def:a07:btorus} are what Ba{\l}aban uses of the rectangular geometry: that the large blocks tile the torus at every scale (the partition of unity and the block operators), and that a cube of two large blocks' side embeds (the local parametrix and the local gauge). A rectangular torus of sides that are multiples of $M_b$ and at least $3M_b$ satisfies both. (ii) The systole condition can be removed by passing to the covering torus $\eta\Z^3/3\Gamma$ and pulling back, at the price of the factor $3^3$ in the kernel bounds from the image sum over the deck group; it is not needed for the tori of the rotation part, whose systole is $10MJ$ at least.
\end{remark}


